%!TEX root = ../../main.tex
\section{한계}
\label{sec:disc-limitations}

다음은 결과를 얻은 범위의 조건이다.

\begin{enumerate}[leftmargin=2em]
  \item \textbf{탐색적 지위.} 유연한 기준선, 점수 분해, 외부 평가, 코호트 분리의 결과는 시험 패치 15.16이 초기 분석에서 노출된 뒤에 생산되었고, 보완 실험은 그 뒤의 진단이다. 구간은 고정된 평가기와 선정 아래의 경기 표본 변동을 기술한다.
  \item \textbf{모델 정의 라벨.} 방향 라벨은 동결 평가기의 변화 부호이며, 관측된 절대적 교전 가치는 없다. 라벨의 일치성과 예측 비교의 안정성은 평가기·종점·표적 정의에 따라 다르고, 작은 변화량 범위에서 가장 민감하다(제 \ref{sec:disc-evaluator} 절).
  \item \textbf{정의의 사람 검증 부재.} 독립적 경계 주석이 없으므로 결과는 검출된 교전에 관한 것이다. 킬 없는 전투는 이 텔레메트리에서 분석 대상 밖에 있다.
  \item \textbf{코호트 범위.} pick(한쪽 참여 1명 이하, 15.16 교전의 약 \sharePickTest\%)은 사전에 제외되었다. 두 코호트는 크기, 양성 비율, 균형 구간 비율이 다르므로 각 결과는 코호트 안에서 읽는다.
  \item \textbf{보정기.} 소규모 교전에서 두 단계 선택 규칙은 모든 모델에 시그모이드를 골랐을 것이며, 계약대로의 항등을 주로, 시그모이드를 민감도로 보고하였고 부호는 같다. 합동 대비의 결론은 이 선택에 의존한다.
  \item \textbf{외부 평가.} 점수 전용이며 어댑터가 없다. 주 두 표본의 구간은 보완 부트스트랩에서 계산하였고, 가족 보정 수준에서는 T NA1을 제외한 세 대비가 불확실하다. 외부 예측 시점에서의 평가기 Brier는 평가기의 품질에 관한 것이며 외부 라벨의 타당성은 별도의 질문이다.
  \item \textbf{정의 상수 민감도.} 사례 구성의 일변량 민감도는 5{,}000 경기의 예산 표본에서 측정하였고, 바뀐 정의에서 동결 $q$·PT의 점수 민감도와 재학습 모드는 미실행이다. 보조 상수의 민감도는 이전 상수 또는 파일럿 코퍼스에서 측정된 값이다.
  \item \textbf{학습기와 정보 구성.} 정보군 비교는 LightGBM 하나의 절차 아래의 결과이고, 동일 352열 입력에서 로지스틱과 LightGBM을 비교하는 정합 실험은 미실행이다. T와 S의 효과 차이는 직접 검정하지 않았다.
  \item \textbf{무결성 검사.} 동결 아티팩트의 digest와 헤드라인 재집계, fold 평가기 재적재는 통과하였다. fold 홀드아웃 경기 집합의 해시 검증과 사후 사건 제거 후의 특징 재구성 불변성은 미완료이고 실행 출처 기록과 번들 재적재 동일성은 부분 완료이므로, 원자료부터 전체 파이프라인의 재현성 검증은 아직 완료되지 않았다. 상태 표현을 만드는 모듈은 해시로 식별되고 열 이름과 블록 구성은 매니페스트에 기록되어 있으며, 참가자 상태 90열의 정규화 상수는 그 모듈의 소스에서 읽어야 한다.
  \item \textbf{미노출 표본.} 미노출 패치의 확인적 평가는 수집 규칙을 잠근 뒤 API 접근 실패로 차단되었다. 본 논문의 외부 결과는 이미 노출된 16.x 표본에 대한 점수 전용 평가이다.
  \item \textbf{다중성.} 주 결과의 열네 개 안팎의 부트스트랩 구간은 보정 없이 보고하였고, 외부 보완 부트스트랩과 정보군 대비만 가족 단위로 등록하였다.
\end{enumerate}
